\section{DDPM 采样算法}

\subsection{反向过程采样}

训练完成后，生成新样本的过程（反向过程 / 采样过程）如下：

\begin{enumerate}
    \item 从标准高斯分布采样初始噪声：$x_T \sim \mathcal{N}(0, I)$
    \item 对 $t = T, T-1, \dots, 1$：
    \begin{enumerate}
        \item 用网络预测噪声：$\hat{\epsilon} = \epsilon_\theta(x_t, t)$
        \item 计算均值：$\mu_\theta(x_t, t) = \frac{1}{\sqrt{\alpha_t}}\left(x_t - \frac{\beta_t}{\sqrt{1-\bar{\alpha}_t}}\, \hat{\epsilon}\right)$
        \item 若 $t > 1$，采样 $z \sim \mathcal{N}(0, I)$；若 $t = 1$，令 $z = 0$
        \item $x_{t-1} = \mu_\theta(x_t, t) + \sigma_t\, z$
    \end{enumerate}
    \item 返回 $x_0$
\end{enumerate}

\begin{importantbox}{采样的直觉：逐步去噪}
采样过程可以理解为：从一团纯噪声开始，每一步都让网络「看一看」当前图像，预测其中有多少噪声，然后减去预测的噪声（同时加入少量新噪声以保持随机性），如此迭代直到获得干净图像。最后一步（$t=1$）不加噪声，因为我们要得到确定性的输出。
\end{importantbox}

\subsection{采样的伪代码}

\begin{lstlisting}[caption={DDPM 采样算法}, label=lst:sample]
def ddpm_sample(model, T, alpha, alpha_bar, beta):
    x = torch.randn(shape)  # x_T ~ N(0, I)

    for t in range(T, 0, -1):
        epsilon_pred = model(x, t)

        # Compute mean
        mu = (1 / alpha[t].sqrt()) * (
            x - (beta[t] / (1 - alpha_bar[t]).sqrt()) * epsilon_pred
        )

        if t > 1:
            sigma = beta[t].sqrt()  # or tilde_beta[t].sqrt()
            z = torch.randn_like(x)
            x = mu + sigma * z
        else:
            x = mu  # No noise at final step

    return x
\end{lstlisting}

\begin{warningbox}{采样速度是 DDPM 的主要瓶颈}
DDPM 的采样需要迭代 $T$ 步（通常 $T = 1000$），每一步都需要一次完整的神经网络前向传播。这意味着生成一张图像需要 1000 次前向传播，相比 GAN 的单次前向传播慢了约三个数量级。后续的 DDIM、DPM-Solver 等工作致力于减少采样步数，可在 10--50 步内获得接近的质量。
\end{warningbox}

\subsection{采样过程中的 $\hat{x}_0$ 可视化}

当使用噪声预测网络时，在每个中间步骤，我们可以利用关系式计算当前对 $x_0$ 的估计：

$$\hat{x}_0^{(t)} = \frac{x_t - \sqrt{1-\bar{\alpha}_t}\, \epsilon_\theta(x_t, t)}{\sqrt{\bar{\alpha}_t}}$$

在采样过程的早期（$t$ 接近 $T$），$\hat{x}_0^{(t)}$ 非常模糊，随着 $t$ 减小，预测逐渐变得清晰。这种可视化直观地展示了扩散模型的``迭代精化''特性。

\begin{knowledgebox}{迭代精化与 Progressive Refinement}
每一步的 $\hat{x}_0$ 预测可以理解为：在当前信息水平下，对最终结果的最佳猜测。随着去噪的推进，可用信息越来越多（噪声越来越少），预测自然越来越准确。这种 progressive refinement 的特性在实际应用中非常有用——例如可以在中间步骤引导生成方向（classifier guidance / classifier-free guidance）。
\end{knowledgebox}

\subsection{本章小结}

DDPM 的采样过程是训练过程的``逆运算''：从纯噪声出发，逐步去噪直到获得干净图像。每一步预测噪声、计算均值、加入受控随机噪声。采样速度慢是主要局限，但后续有大量加速方法。

%% ========== 第八节：扩散模型的表达能力 ==========

